\subsection{Cascading}
\label{lit:cascading}
Modern RAG systems encounter diverse workloads, ranging from simple informational queries to highly complex analytical tasks. In production-grade enterprise environments — where high availability, low latency, and cost efficiency are critical constraints — a "one-size-fits-all" architecture is fundamentally sub-optimal \cite{jeong2024adaptive, chen2023frugalgpt}. It over-allocates expensive computational resources for trivial questions, yet may lack the retrieval depth and reasoning capacity required for complex queries. To mitigate this, model cascading or query routing can be employed to dynamically direct inputs to an appropriately scaled configuration of the RAG pipeline \cite{ong2025routellm}. For instance, a simple query can bypass the full pipeline and utilize a lightweight LLM with minimal context. Conversely, a complex query requiring synthesis across multiple enterprise knowledge bases can trigger an augmented pipeline with extended retrieval top-k parameters and a more capable foundation model. 

Model cascading offers substantial operational cost reductions both in isolated Large Language Model (LLM) invocations and in retrieval-augmented generation (RAG) pipelines whose overall cost profiles are dominated by API usage. For instance, FrugalGPT \cite{chen2023frugalgpt} routes individual requests sequentially from lightweight to heavyweight models based on a reliability score generated by a scoring classifier. This cascading approach achieves up to a 98.3\% cost reduction compared to a homogeneous GPT-4 baseline while maintaining comparable response quality. Similarly, RouteLLM \cite{ong2025routellm} yields a 1.14$\times$ to 3.66$\times$ cost reduction compared to a static GPT-4 deployment by dynamically bypassing expensive models, preserving approximately 95\% of the baseline's quality. Similar cost reduction patterns have also been demonstrated by other adaptive RAG frameworks \cite{jeong2024adaptive, guo2026deepsieve}.

Routing low-complexity queries to smaller models yields significant latency benefits due to their faster inference speeds. Crucially, the computational overhead introduced by the routing decision itself remains minimal, typically amounting to approximately 11\, ms per query \cite{zhang2025ragrouter} or a 0.4\% latency increase \cite{ong2025routellm}. In straightforward scenarios, bypassing retrieval operations entirely can yield order-of-magnitude improvements: \citet{jeong2024adaptive} reports a reduction in generation latency from 3.08 and 27.18 seconds (for single-step and multi-step retrieval pipelines) to 0.35 seconds. Furthermore, adaptive stopping mechanisms, such as that proposed by \citet{park2025stoprag}, can prune unnecessary retrieval-generation loops, reducing the total number of iterations by 14--49\%. 

While providing both operational cost and latency benefits, query routing has a modest impact on quality. For example, \citet{jeong2024adaptive} found that the F1 score remained practically unchanged, while achieving a substantial reduction in latency (~56\%) and in token consumption (63\%) compared to a multi-step retrieval baseline (using GPT-3.5). Similar results have been shown by \citet{guo2026deepsieve}, where token consumption decreased by up to ~90\% while losing only 0.8 points on F1, but improving the Exact Match score by 2.6 points compared to the agentic baseline, and yielding 20.7 and 15.5 percentage point improvements in F1 and Exact Match (EM) scores, respectively, compared to naive RAG.

\subsection{Small-to-Big retrieval}
\label{lit:retrieval}

 Document chunking is a foundational phase in RAG pipelines, directly determining the granularity of the vector index. However, selecting a static chunk size introduces a fundamental trade-off between retrieval precision and contextual completeness. While smaller chunks are computationally efficient and exhibit high intrachunk cohesion (ICC), they frequently suffer from low document contextual coherence (DCC) due to semantic fragmentation \cite{demourajunior2026adaptive}. Conversely, larger chunks preserve rich contextual relationships but degrade retrieval accuracy by diluting query-relevant semantic features within dense vector representations. Small-to-big retrieval resolves this tension by decoupling the representation used for vector search from the context provided to the generator. In this architecture, dense retrieval is executed over fine-grained index nodes, while the generation phase utilizes a broader, mapped context (such as parent chunks or entire source documents) retrieved via sibling pointer mapping. Implementations of this technique range from neighboring chunk expansion, which dynamically retrieves adjacent sliding windows to restore structural continuity \cite{lumer2025rethinking}, to hierarchical abstract trees, which map granular leaf nodes to high-level thematic summaries\cite{li2026bridgerag}.

Different chunking and representation methods directly increase the accuracy of RAG pipelines. For example, \citet{lumer2025rethinking} showed that expanding the context to neighboring chunks can increase the win rate to 65\% compared to using single independent chunks. In addition, the adaptive chunking method, where chunks are generated dynamically based on document structure, increases downstream answer correctness by 8--10 percentage points without requiring downstream prompt modifications \cite{demourajunior2026adaptive}. Hierarchical abstract trees like Bridge-RAG \cite{li2026bridgerag} yield up to a 44\% improvement in BLEU scores on complex reasoning tasks by tracing parent-child lineages. Similarly, dynamic reasoning frameworks like LogicRAG \cite{chen2026logicrag} and QPaug \cite{kim2024qpaug} increased string accuracy by 14.7\% absolute (compared to the next best graph baseline) and answer F1 scores by over 34\% relatively.

While small-to-big strategies provide substantial benefits to answer correctness and retrieval recall, they introduce operational and computational overhead. At query time, asynchronous parallel search for neighboring chunks \cite{lumer2025rethinking} keeps latency minimal, adding a negligible $\approx 0.17$\,\text{s} overhead. However, more complex reasoning structures like LogicRAG \cite{chen2026logicrag} more than double the query latency
compared to naive flat retrieval, increasing it from 4.28\,\text{s} to 9.83\,\text{s}. Furthermore, advanced chunking and indexing methods are highly time-consuming during the ingestion phase. The slowest model-driven and layout-aware chunking methods can run for tens of hours, whereas fixed-size chunking completes in under a second \cite{smigielski2026chunking}. In hierarchical structures, generating node summaries and index trees also increases offline preprocessing costs (e.g., scaling from \$1.19 to \$6.12 per average filing in financial domains) \cite{lumer2025rethinking}. Nonetheless, this preprocessing latency is typically acceptable in enterprise environments, as chunking and indexing are offline processes that occur only during initial document ingestion.

\subsection{Semantic Re-ranking}
\label{lit:re_rank}


\subsection{Prompt Caching}
\label{lit:caching}


